\chapter{Python API and reproduction commands}
\label{app:python}

The release code is intentionally compact. Its purpose is to verify the algebra and provide a reusable seed rather than to compete with optimized quantum simulators.

\section{Core API}
\begin{lstlisting}[style=pythonstyle,caption={Core packet API.}]
from symmetry_packets import (
    weight_basis,
    residue_basis,
    symmetric_subspace_dimension,
    boson_fixed_number_dimension,
    fermion_fixed_number_dimension,
    PacketPlan,
    rz,
    xy_exchange,
    hypercube_quotient,
    circuit_leakage_bound,
    hamiltonian_leakage_bound,
)
\end{lstlisting}

\section{Exact packet plan}
\begin{lstlisting}[style=pythonstyle,caption={Fixed-weight packet compilation.}]
plan = PacketPlan.fixed_weight(n=10, k=2)
packet_state = random_packet_state(plan, seed=12)
full_state = plan.embed(packet_state)

restricted_gate = plan.restrict_gate(
    xy_exchange(theta=0.2),
    qubits=(3, 4),
)
packet_state = restricted_gate @ packet_state
\end{lstlisting}

\section{Validation}
From the project root, run
\begin{lstlisting}[style=pythonstyle]
PYTHONPATH=code python3 code/validate.py
\end{lstlisting}
The expected output reports successful formula checks and reconstruction errors near machine precision for the representative fixed-weight circuits.

\section{Figure generation}
All figures in the PDF are generated by
\begin{lstlisting}[style=pythonstyle]
python3 code/make_figures.py
\end{lstlisting}
The generator writes vector PDF figures into the \texttt{figures} directory.

\section{Build}
The complete rebuild command is
\begin{lstlisting}[style=pythonstyle]
./run_all.sh
\end{lstlisting}
It validates the Python code, regenerates figures, compiles the LaTeX source, and updates the checksums.
